\section{Lecture 9 (September 30th)}
\begin{defi}
(14.2) Let $f(t)\in K[t]$. A splitting field for $f(t)$ is a field extension $F/K$ such that $f(t)$ factors into linear terms 
\[c\cdot \prod ^{n}_{i=1}(t-\alpha _{i})\]
with $F=K(\alpha_1,\alpha_2,\ldots ,\alpha _{n})$. 
\end{defi}
\vspace{2ex}
\begin{defi}
(13.19) A field $F$ is  algebraically closed if every nonconstant polynomial $f(t)\in F[t]$ has a zero in $F$.
\end{defi}
\vspace{2ex}
\begin{rmk}
If $K={\bm Q}$, the existence of a the splitting field of $f(t)$ is immediate: if $f(t)=c\cdot \prod^{n}_{i=1}(t-\alpha _{i})$ in ${\bm C}[t]$, then
\[F=K(\alpha_1,\ldots ,\alpha _{n})\subset {\bm C}\]
\end{rmk}
\vspace{2ex}
\begin{thm}
(14.6) There exists a splitting field of $f(t)\in K[t]$. 
\end{thm}
\vspace{2ex}
\begin{proof}
We prove this theorem inductively. Assume $\mathrm{deg}(f)=1$. We simply take $L=K$ and $L/K$ is a field extension such that $f(t)$ factors into linear terms. 

\bigskip

Assume that the statement holds for all polynomials with $\mathrm{deg}(f)< n$. For a polynomial $f(t)$ of degree $n$, we can pull out an irreducible factor $p(t)$ and using Kronecker's theorem, we have a field $K_1$ such that $f(t)=(t-\alpha )^{m}f_1(t)$ in $K_1$ with $\mathrm{deg}\,f_1(t)< n$. By assumption, we can find a splitting field of $K_1$ where $f_1(t)$ splits into linear factors:
\[f(t)=(t-\alpha )^{m}\cdot c\cdot \prod^{n-1}_{i=1}(t-\alpha _{i})\in L\]
and $L/K$ is a field extension such that $f(t)$ factors into linear terms with $L=K(\alpha_1,\alpha_2,\ldots ,\alpha_{n})$.
\end{proof}
\vspace{2ex}
\begin{thm}
(14.3) Let $K$ be a field. Then there exists an algebraic extension $\bar{K}/K$ such that $\bar{K}$ is algebraically closed.
\begin{center}
\begin{tikzpicture}[baseline=(current bounding box.center)]
  \node (F) at (0,3) {$\bar{K}$};
  \node (K) at (0,0) {$K$};
  \node (L) at (2,1.5) {$L$};

  % lines (like the chalk sketch)
  \draw (F) -- (L);
  \draw (K) -- (L);
  \draw (F) -- (K);
\end{tikzpicture}
\end{center}
\end{thm}
\vspace{2ex}
\begin{rmk}
Such an extension is unique: if $F$ and $L$ are two such extensions of $K$, then there exists an isomorphism $\sigma :F\rightarrow L$ such that $\sigma |_{K}=\mathrm{id}_{K}$ (that is, $\sigma (a)=a$ for all $a\in K$).
\begin{center}
\begin{tikzpicture}[baseline=(current bounding box.center),
  >={Straight Barb[length=3pt,width=6pt]}, % barbed arrowhead (side lines)
  node font=\small]

  % corners
  \node (A) at (0,3) {$F$};
  \node (B) at (0,0) {$K$};
  \node (C) at (3,0) {$K$};
  \node (D) at (3,3) {$L$};

  % vertical arrows (pointing up)
  \draw[->] (B) -- (A);
  \draw[->] (C) -- (D);

  % horizontal arrows (pointing right) with labels above
  \draw[->] (A) -- node[midway, above] {$\sigma $} (D);
  \draw[->] (B) -- node[midway, above] {${\bf 1}$} (C);

  % commutator symbol in the middle (clockwise)
  \node at (1.5,1.5) {\scalebox{2}{$\circlearrowright$}};

\end{tikzpicture}
\end{center}

\end{rmk}
\begin{defi}
(14.4) We call $\bar{K}$ the algebraic closure of $K$.
\end{defi}
\vspace{2ex}




